% weyl_energy.tex -- round 389, lemma-first attack on the Weyl
% energy of cancellation after four majorant routes died.
% Outcome: REDUCED.  Builds with pdflatex.
% Companion machine check: verify_weyl_energy.py.
% Discovery census: experiments/tfpt-discovery/weyl_energy_probe.py.
% NO RH CLAIM.  Research documentation, not a theorem of RH.
\documentclass[11pt]{article}

\usepackage[a4paper,margin=2.7cm]{geometry}
\usepackage{amsmath,amssymb,amsthm}
\usepackage{booktabs}
\usepackage{microtype}
\usepackage{xcolor}
\usepackage[colorlinks=true,linkcolor=blue!50!black,urlcolor=blue!50!black]{hyperref}

\newtheorem{lemma}{Lemma}
\newtheorem{prop}{Proposition}
\newtheorem{definition}{Definition}
\theoremstyle{remark}
\newtheorem{remark}{Remark}

\newcommand{\R}{\mathbb{R}}
\newcommand{\Q}{\mathbb{Q}}
\newcommand{\N}{\mathbb{N}}
\newcommand{\eps}{\varepsilon}

\title{\bfseries The Weyl energy of cancellation, reduced\\[2mm]
\large Exact Fourier decomposition of the three open objects,
Parseval on the cosine grid, and the scale that actually closes}
\author{Stefan Hamann}
\date{August 28, 2026}

\begin{document}
\maketitle

\begin{abstract}\noindent
Four majorant routes died at the same place
(r381 $L^\infty$, r385 Koksma, r387 Gershgorin, r388 osc-Geronimus).
The three open cancellation objects are the signed $\mu$-CD
off-diagonal (assist $0.0399$ against Gershgorin $13.32$),
terminal $C_\varepsilon$ ($\lvert\mathrm{FO}\rvert\le\tfrac1{32}$
on the last twelve), and $Z_{\mathrm{loc}}$ (median cancel $0.055$,
not uniform).
This note replaces majorization by an exact Fourier identity.
What is proved: on the cosine grid the Chebyshev-$T$ Christoffel--Darboux
kernel is a finite cosine sum, and
\[
\Phi_k
:=
\sum_{i,j} v_i v_j K_k^T(x_i,x_j)
=
\frac1\pi\lvert S_0\rvert^2
+
\frac2\pi\sum_{m=1}^{k-1} C_m^2,
\]
with $S_q=\sum_j v_j e^{iq\theta_j}$ and $C_m=\mathrm{Re}\,S_m$;
the off-diagonal is $\Phi_k$ minus an explicit $v^2$-Weyl diagonal;
$Q_k^T=(C_0+C_{2k})_\nu/(C_0+C_{2k})_\mu$;
$\pi_k^2$ is a finite combination of $C_0,\ldots,C_{2k}$ via the
Chebyshev coefficients of $\pi_k$;
and the finite-grid Parseval
$\sum_{q=0}^{2S-1}\lvert S_q\rvert^2=2S\sum_j v_j^2$
holds with ratio $1$ on the two-period mesh and on FRAME-A $w=9$.
What is \emph{not} proved: a bound on $\sum c_q\lvert S_q\rvert^2$
that closes all three objects.
MAIN and $\chi_3$ sit at quadratic-mean
($\mathrm{mean}_{m\ge1} C_m^2\big/({\textstyle\frac12}\sum v^2)\in[0.66,1.02]$
on core-$42$);
the two-period is a spectral comb ($\rho_{\mathrm{AP}}=1$ if and only
if the energy concentrates on a progression);
scramble holds quadratic-mean and dies on the $\mu$-OP remainder.
The Chebyshev piece of $C_\varepsilon$ closes at quadratic-mean
($\mathrm{FO}^T\le0.0165<\tfrac1{32}$ on the eleven small-$N$ windows).
Assist is not the Chebyshev all-ones energy
($1^TE^T1/n_\nu=0.114$ against $\lambda=0.99983$).
$Z_{\mathrm{loc}}$ is not the Chebyshev $x T_k^2$ proxy (cancel $0.47$).
The hydra-legal remainder is the Weyl energy of the $\mu$-OP
Chebyshev-coefficient Gram: of $\phi^*$ (assist), of consecutive
$\alpha_k$ ($d\Delta$/ $C_\varepsilon$), of the edge-masked drive
($Z_{\mathrm{loc}}$).
No RH claim.
\end{abstract}

\medskip\noindent
\textbf{Status.}\enspace
Lemmas~\ref{lem:energy}--\ref{lem:parseval} are proved (finite identities).
Lemma~\ref{lem:QT} and Lemma~\ref{lem:pik} are proved (finite identities).
Proposition~\ref{prop:spectrum} is a finite machine census.
Proposition~\ref{prop:scale} records the scale adjudication.
Proposition~\ref{prop:red} names the reduction.
No RH claim.  No $L^*$ claim.  No $C_\varepsilon$ claim.

\section{The object}\label{sec:obj}

Write $\theta=\arccos x$ on the cosine grid of~\cite{r387},
$S_q=\sum_j v_j e^{iq\theta_j}$ for the $\nu$-weights, and
$C_m=\sum_j v_j\cos(m\theta_j)=\mathrm{Re}\,S_m$.
The Chebyshev-$T$ CD kernel of $k$ orthonormal functions is
\begin{equation}\label{eq:K}
K_k^T(\cos\theta,\cos\phi)
=
\frac1\pi
+
\frac2\pi\sum_{m=1}^{k-1}\cos(m\theta)\cos(m\phi)
=
\frac{D_{k-1}(\theta-\phi)+D_{k-1}(\theta+\phi)}{2\pi},
\end{equation}
the r387 identity.
The IIKS Gram is $(E_k)_{ij}=\sqrt{v_i v_j}\,K_k^\mu(y_i,y_j)$
with $K_k^\mu$ the $\mu$-CD kernel;
r285 writes $\lambda_{\max}(E_k)=\mathrm{maxdiag}_k\cdot(1+\mathrm{assist}_k)$.
Round 385: $\mathrm{FO}_k[-\nu]=\gamma_k(Q_k-Q_{k-1})$ with
$Q_k=Q_k^T+\Delta_k$.
$Z_{\mathrm{loc}}$ is the edge piece of the r383 canonical split.

The candidate of this round: compute the three cancellations as
explicit linear combinations of $C_m$ (or $\lvert S_q\rvert^2$),
then read off the scale of a closing bound from the measured
spectrum.

\section{What is exact}\label{sec:exact}

\begin{lemma}[Chebyshev energy]\label{lem:energy}
For any atomic weights $v$ on $(-1,1)$,
\[
\Phi_k
=
\sum_{i,j}v_i v_j K_k^T(x_i,x_j)
=
\frac1\pi C_0^2
+
\frac2\pi\sum_{m=1}^{k-1} C_m^2.
\]
The off-diagonal is $\Phi_k-\sum_i v_i^2 K_k^T(x_i,x_i)$, with
\[
K_k^T(x,x)
=
\frac k\pi
+
\frac1\pi\sum_{m=1}^{k-1}\cos(2m\theta),\qquad
\sum_i v_i^2 K_k^T(x_i,x_i)
=
\frac k\pi\sum_i v_i^2
+
\frac1\pi\sum_{m=1}^{k-1} W_{2m},
\]
$W_q=\sum_i v_i^2\cos(q\theta_i)$.
Exact over $\Q$ on a five-atom toy ($\Phi^*=99$, off $=26$,
$\mathrm{diag}=73$, the $1/\pi$ scale stripped).
Floating-point six-atom residual $0$ (energy) and $3.4\cdot10^{-16}$ (off).
FRAME-A $w=9$, $k=8$: residual $2.6\cdot10^{-15}$;
mpmath dps $30$ residual $4.3\cdot10^{-31}$.
\end{lemma}

The $\mu$-CD remainder is carried explicitly:
$K_k^\mu=K_k^T+R_k$, so the $v\otimes v$ bilinear splits as
$\Phi_k^\mu=\Phi_k+\sum_{i,j}v_i v_j R_k$.
At $k=22$ on $w=9$ this remainder is $-0.0028$ against
$\Phi=0.026$ (small).
The $\sqrt{v}$-dressed Gram $E$ is a \emph{different} quadratic form:
$1^TE1=16.05$ against $1^TE^T1=6.44$ at the same $k$
(remainder of order one).

\begin{lemma}[$Q^T$ as a Weyl ratio]\label{lem:QT}
\[
Q_k^T
=
\frac{C_0(\nu)+C_{2k}(\nu)}{C_0(\mu)+C_{2k}(\mu)}.
\]
This is~\cite{r385}'s Chebyshev sampling, rewritten in $C_m$.
Checked on a six-atom toy at $k=0,1,2,4$ (residual $0$) and on
FRAME-A $w=9$ at $k=0,10,50,182$.
\end{lemma}

\begin{lemma}[$\pi_k^2$ Weyl]\label{lem:pik}
If $\pi_k=\sum_{m=0}^k\alpha_m T_m$ in the Chebyshev basis, then
\[
Q_k
=
\sum_j v_j\pi_k(x_j)^2
=
\sum_{p,q=0}^k\alpha_p\alpha_q\cdot\tfrac12\bigl(C_{p+q}+C_{\lvert p-q\rvert}\bigr).
\]
Five-atom toy, degree $3$, residual $5\cdot10^{-16}$.
Consequently $\mathrm{FO}_k=\gamma_k(Q_k-Q_{k-1})$ is a linear
combination of $C_0,\ldots,C_{2k}$, plus the r388 split into
$dQ^T+d\Delta$.
\end{lemma}

\begin{lemma}[finite-grid Parseval]\label{lem:parseval}
On the cosine index grid $\theta_j=\pi j/S$, $j=1,\ldots,S$,
for any (possibly sparse) real weights $v$,
\[
\sum_{q=0}^{2S-1}\lvert S_q\rvert^2
=
2S\sum_j v_j^2.
\]
Ratio $=1$ at machine precision on the two-period mesh $S=21$
and on FRAME-A $w=9$ (nu-only and the kept union of $367$ nodes).
\end{lemma}

This is the exact Parseval balance asked of the round.
It does \emph{not} by itself imply quadratic-mean per mode:
the two-period saturates Parseval \emph{and} concentrates.

\section{The spectrum}\label{sec:spec}

Write $\mathrm{qm}:=\tfrac12\sum v^2$ for the quadratic-mean
reference ($E[C_m^2]$ if a mode were a random phase on the circle)
and $\mathrm{HHI}(C)=\sum_{m=1}^{M} p_m^2$ for the energy share
$p_m\propto C_m^2$.

\begin{prop}[census]\label{prop:spectrum}
On the sealed constructors:
\begin{enumerate}
\item FRAME-A $w=9$, $M=400$:
$\mathrm{mean}/{\mathrm{qm}}=0.922$,
$\max/{\mathrm{qm}}=14.5$ at $m=208$,
$\lvert C\rvert/\mathrm{mass}$ last-twelve in $[0.0005,0.179]$
(trivial Weyl, not PNT),
$\mathrm{HHI}_{1..80}=0.053$ (random $\sim0.013$).
\item Core-$42$, $M=200$:
$\mathrm{mean}/{\mathrm{qm}}\in[0.661,1.018]$,
$\lvert C\rvert/\mathrm{mass}$ max in $[0.229,0.401]<\tfrac12$,
$\mathrm{HHI}\in[0.013,0.097]<0.12$.
\item EXT-heavy seven:
$\lvert C\rvert/\mathrm{mass}$ max in $[0.237,0.348]<0.40$;
$\mathrm{HHI}\in[0.16,0.35]$ is flank-mass, not a two-period comb.
\item $\chi_3$-$42$: $\mathrm{mean}/{\mathrm{qm}}\in[0.658,0.909]$,
same class as MAIN.
\item Two-period $S=81$, $c=\tfrac23$:
$\lvert C_S\rvert/\mathrm{mass}=1$,
$\mathrm{HHI}_{1..80}=0.727$,
comb at $q=0,1,S-1,S,S+1$.
\item Scramble seed $1$: $\mathrm{mean}/{\mathrm{qm}}=0.947$
(quadratic-mean holds), $\mathrm{HHI}=0.007$ (more uniform than MAIN).
Top $\lvert\Delta C_m^2\rvert$ modes $380,354,30,177,170,23$.
\item Von Mangoldt versus equal weight on $w=9$, $M=200$:
weighted $\mathrm{mean}/{\mathrm{qm}}=0.718>$ equal $0.403$.
The weights do not buy extra cosine cancellation.
\end{enumerate}
\end{prop}

The cancellation of the program, at the Chebyshev layer, sits at
quadratic-mean: no structured valley below $\mathrm{qm}$, no
resonance at a fold frequency on MAIN.
The two-period is the concentrating adversary, and $\rho_{\mathrm{AP}}=1$
is exactly that comb.
Scramble does not die on $\nu$-Weyl concentration.

\section{Scale adjudication}\label{sec:scale}

\begin{prop}[scale]\label{prop:scale}
\begin{enumerate}
\item[(a)] \emph{Assist.}
Need $\mathrm{assist}<(1-\mathrm{maxdiag})/\mathrm{maxdiag}=0.0401$
at the $w=9$ wall (true $0.0399$).
The Chebyshev all-ones Rayleigh is $1^TE^T1/n_\nu=0.114\ll\lambda=0.99983$;
the $v\otimes v$ Chebyshev off-diagonal at the wall is \emph{negative}
($-0.007$).
A quadratic-mean majorant of $1^TE^T1$ overshoots the actual
Chebyshev energy by $1.52$.
Quadratic-mean of the Chebyshev bilinear does \emph{not} close assist:
it is the wrong quadratic form (all-ones, not the top mode of $E$).
\item[(b)] \emph{$C_\varepsilon$.}
Need $\lvert\mathrm{FO}\rvert\le\tfrac1{32}$ on the last twelve.
$\mathrm{FO}^T=0.00496<\tfrac1{32}$ on $w=9$;
eleven small-$N$ windows give $\max\lvert\mathrm{FO}^T\rvert=0.01650$
and $\max\lvert\mathrm{FO}\rvert=0.01182$, both below the bar.
The Chebyshev piece closes at quadratic-mean
($\lvert C_{2k}\rvert/\mathrm{mass}$ last-twelve mean $0.089$ against
rms-prediction $0.105$).
The live remainder is $d\Delta$ (r388).
\item[(c)] \emph{$Z_{\mathrm{loc}}$.}
The Chebyshev proxy $\sum v\cos\theta\,(1+\cos 2k\theta)/2$
has signed/abs cancel $0.475$, against r386's median $0.055$.
The edge restriction $\lvert x\rvert>0.85$ is \emph{worse} ($0.86$).
$Z_{\mathrm{loc}}$ is the edge-masked drive of the actual $\mu$-OP,
not this proxy; quadratic-mean of the proxy does not close uniform cancel.
\end{enumerate}
\end{prop}

PNT never enters: $\lvert C\rvert/\mathrm{mass}\le0.40$ is already
the trivial Weyl size, and the Chebyshev $\mathrm{FO}^T$ bound
does not demand anything smaller.

\section{Reduction}\label{sec:red}

\begin{prop}[reduced]\label{prop:red}
The energy identity and the finite-grid Parseval are theorems.
MAIN$/\chi$ are at quadratic-mean; $\rho_{\mathrm{AP}}<1/5$ is the
time-domain form of ``not a comb''.
The deduction ``therefore assist, $C_\varepsilon$ and $Z_{\mathrm{loc}}$
close at quadratic-mean'' is false for assist and $Z_{\mathrm{loc}}$,
and true for the Chebyshev piece of $C_\varepsilon$.
The hydra-legal remainder is the Weyl energy of the $\mu$-OP
Chebyshev-coefficient Gram: of the top eigenvector $\phi^*$ of $E$
(assist), of consecutive Chebyshev coefficients $\alpha_k$ of $\pi_k$
($d\Delta$, the live $C_\varepsilon$), and of the edge-masked drive
($Z_{\mathrm{loc}}$).
\end{prop}

\section{Kills}\label{sec:kills}

Wrong $c_q$ ($1$ instead of $2$) fails at relative error $0.108$.
Dropping the diagonal fails at $\lvert\Phi-\mathrm{off}\rvert/\lvert\mathrm{off}\rvert=2.78$.
The two-period is visible as a $q$-comb ($\lvert C_S\rvert/\mathrm{mass}=1$,
$\mathrm{HHI}=0.727$); this is $\rho_{\mathrm{AP}}$ as a spectral statement.
Scramble seed $1$ holds quadratic-mean and is \emph{more} uniform than MAIN;
the divergence modes of $\lvert\Delta C_m^2\rvert$ are
$380,354,30,177,170,23$.
It dies on the $\mu$-OP remainder ($d\Delta=2.16\cdot10^6$ of~\cite{r388}),
not on $\nu$-Weyl concentration.

\begin{thebibliography}{9}
\bibitem{r381} S.~Hamann, \emph{The $G_\varepsilon$ lemma, reduced},
rh/problem/g\_eps\_lemma.tex, 2026.
\bibitem{r385} S.~Hamann, \emph{Christoffel quietness, reduced},
rh/problem/christoffel\_quiet.tex, 2026.
\bibitem{r386} S.~Hamann, \emph{Living-ladder $(Z')$ and cofinal $R_0$},
rh/problem/compose\_premises2.tex, 2026.
\bibitem{r387} S.~Hamann, \emph{Coherence assist, reduced},
rh/problem/coherence\_assist.tex, 2026.
\bibitem{r388} S.~Hamann, \emph{The $\Delta$-deformation, reduced},
rh/problem/delta\_deformation.tex, 2026.
\end{thebibliography}

\end{document}
